\documentclass[11pt,a4paper]{article}
\usepackage[margin=23mm]{geometry}
\usepackage{amsmath,amssymb,booktabs,longtable,array,hyperref}
\hypersetup{hidelinks}
\emergencystretch=3em
\newcommand{\C}{\textbf{(c)}\nobreakspace}
\newcommand{\D}{\textbf{(d)}\nobreakspace}
\newcommand{\dd}{\mathrm d}
\newcommand{\file}[1]{\nolinkurl{#1}}
\newcommand{\checks}[1]{\par\noindent\textit{Checks:} \file{twobody_toda_#1.py}.\par}
\newcolumntype{L}[1]{>{\raggedright\arraybackslash}p{#1}}
\title{The Toda two-body edge phase at fixed dressed momenta:\\
directional one-loop coefficient, exchange decomposition,\\
and the surviving physical-normalization question}
\author{Research note for Gustav Delius}
\date{4 November 2026}
\begin{document}
\maketitle
\begin{abstract}
\C Conditional on the corrected half-line asymptotics, the inherited
\emph{two-term directional} Toda phase has the following leading
small-amplitude, fixed-dressed-momentum correction. For right-moving
overtaking solitons, $r=\beta_2/\beta_1\in(0,1)$,
\begin{align*}
 R_{12}^{(1)}&=-\frac4\pi[\chi_2(r)-r],\\
 S^{(1)}&=\frac2\pi(r+r^{-1})-\pi,\\
 \vartheta_N^{(1)}&=\pi-\frac4\pi\chi_2(r)+\frac2\pi(r-r^{-1}),
\end{align*}
where $S=(R_{12}+R_{21})/2$, $\vartheta_N=(R_{12}-R_{21})/2$, and
$\chi_2(r)=[\operatorname{Li}_2(r)-\operatorname{Li}_2(-r)]/2$.
The reverse ordered branch is supplied by exchange, not by substituting
$r>1$ into principal real logarithms. The corresponding leading
opposite-edge directional coefficients vanish. New finite-sea solves,
high-precision inversions and unconstrained intercept fits test these
statements; the numerical budgets are given below.

\C The same momentum logarithm fixes the symmetric term and the one-body
invariant $C-s_1=-1/(4\pi)$. We also find a two-term/four-term mismatch
in the Krein equation printed in the frame audit. Its correction does not
restore an arbitrary antisymmetric Toda freedom.
\D These results do \emph{not} identify $\vartheta_N$ with the
Gamma-relative continuum residual, nor establish a uniform fixed-nonzero-momentum
semiclassical limit. The exact endpoint-bilinear representative has additional
one-body reference terms, including half-integer powers in the edge expansion.
They are explicitly retained. Thus the literal fully normalized physical
scalar phase requested in the task is not established by this package.
\end{abstract}

\section{Scope, conventions and the observable actually computed}
Labels retain the repository meanings: (a) proved in literature, (b)
literature heuristic, (c) project derivation or numerical evidence, (d)
unestablished interpretation. No external paper, abstract, arXiv entry or DOI
is cited or claimed newly read. Internal sources are listed at the end.
All new scripts and caches have the required \file{twobody_toda} prefix.
No earlier script is run or imported, and no git state is touched.

\C Retain the audited units $a=1$, $b=e^{-1/2}$,
\begin{equation}
 h=\hbar_T=3b\hbar_B,\qquad
 \kappa=\sqrt3\beta=\sqrt{hy/b},\qquad y=\beta^2/\hbar_B,
 \qquad P_B=\overline P/(6\sqrt3b).
 \label{eq:units}
\end{equation}
The inherited continuum is the bad equation
$u_{tt}=u_{xx}+u_{xxxx}+3(u^2)_{xx}$, $u=2(\log\tau)_{xx}$,
$k=\sinh\beta$, $v=\epsilon\cosh\beta$.
The dictionary is a small-amplitude dictionary; no finite-amplitude
Toda--Boussinesq equivalence is inferred.
In physical Boussinesq momenta the leading ratio is
$r=(|P_{B,2}|/|P_{B,1}|)^{1/3}$, since
$|P_B|=2\beta^3/3+O(\beta^5)$ in this corner.
The exact classical Toda quantities used in the numerical subtractions are
\begin{align}
 P_0(\kappa,\epsilon)&=4\epsilon b(\kappa\cosh\kappa-\sinh\kappa),
 &p_0&=2\epsilon b\cosh\kappa,\\
 M_0&=-2\kappa,&\Pi_0&=2\epsilon b\sinh\kappa,
 &v_0&=\epsilon b\sinh\kappa/\kappa.
\end{align}
At each finite $h$ we solve $\overline P(p_j)=P_0(\kappa_j,\epsilon_j)$;
we never replace the rapidity coordinate by a dressed momentum.

\C Set $e_j=\epsilon_jq$ and define the directional phase and its
four-term counterpart by
\begin{align}
 N_{12}&=\Phi(p_1,p_2)-\Phi(p_1,e_2),\\
 U_{12}&=\Phi(e_1,p_2)-\Phi(e_1,e_2),&
 \widehat\Phi_{12}&=N_{12}-U_{12}.
 \label{eq:references}
\end{align}
Our residual is precisely the audited Phase 2f/2h two-term quantity,
with the \emph{full classical lattice} subtraction:
\begin{equation}
 R_{12,h}=\frac{N_{12}-\delta_T(1,2)-C_{12}}h,
 \qquad C_{12}=M_{0,1}(P_{0,2}-\Pi_{0,2}).
 \label{eq:R}
\end{equation}
Here $\delta_T=(p_{0,1}-p_{0,2})\log A_T+M_{0,1}M_{0,2}(v_{0,1}-v_{0,2})$,
with $A_T=\sinh^2[(\kappa_1-\kappa_2)/2]/\sinh^2[(\kappa_1+\kappa_2)/2]$
for equal directions and the same ratio of $\cosh^2$ for opposite directions.
The subtraction is evaluated with 60-digit arithmetic before comparison to
the float64 sea solves. The dimensionless action phase is unchanged by the
audited action rescaling: $\delta_T/h=\delta_B/\hbar_B$ in the scaling limit.

\D A directional thermodynamic Bethe phase need not be antisymmetric.
Consequently the notation $\vartheta_N$ is deliberate: it denotes the
exchange-odd part of \eqref{eq:R}, not an already normalized scalar
S-matrix, and not the residual relative to a Gamma candidate in
\file{smx_candidates.tex}. This distinction matters especially for head-on
scattering and absolute phase constants.
\checks{scan}\checks{identities}

\section{Exact finite-sea identities and the reference term}
\C With $t=2\pi h$,
\begin{align}
 \Theta(d)&=2h\operatorname{Im}\log\Gamma(1+id/h)+2d\log h,
 &K(d)&=\Theta'(d),\\
 (\mathcal Af)(x)&=tf(x)-\int_{-q}^qK(x-z)f(z)\dd z,
 &\mathcal A\rho&=1,\quad \mathcal A\eta=x,\quad\int\rho=1,\\
 \mathcal AF_z(x)&=\Theta(z-x),
 &\Phi(p,z)&=\Theta(p-z)-\int_{-q}^qK(p-x)F_z(x)\dd x.
\end{align}
The gamma logarithm is continuously unwrapped on the real difference axis.
Let $D=P'$, $V=E'$, $D_q=t\rho(q)$, $V_q=t\eta(q)$.
The audited endpoint identities and Jacobian are
\begin{align}
 D_q[F_p(q)-F_p(-q)]&=1-D(p),\\
 V_q[F_p(q)+F_p(-q)]&=P(p)-V(p),\\
 \frac{\dd u}{\dd y}&=\frac{2\sinh\kappa}{D(q+hu)},
 &\overline P(q+hu)&=P_0(\sqrt{hy/b}).
 \label{eq:jac}
\end{align}
The high-precision checker also uses an independently counted spacing
$a\ne1$, restoring $a$ in the Jacobian denominator, and differentiates the
inverted momentum numerically. Its largest $N=48$ Fredholm identity residual
is $2.57\,10^{-30}$; the centered-difference Jacobian error is
$6.65\,10^{-22}$.

\C Evaluating the defining equation for $F_z$ at an endpoint gives the
particularly useful identity
\begin{equation}
 \boxed{\Phi(e,z)=-tF_z(e),\qquad
 U_{12}=-t[F_{p_2}(e_1)-F_{e_2}(e_1)].}
 \label{eq:Uexact}
\end{equation}
This fixes the one-body reference term without a fit or a phase convention.
Writing $f_j^\pm=F_{p_j}(\pm q)-F_{e_j}(\pm q)$, the exact reciprocity law is
\begin{equation}
 \widehat\Phi_{12}+\widehat\Phi_{21}
 =t(f_1^+f_2^+-f_1^-f_2^-)
 =\mathfrak M_1\mathfrak Q_2+\mathfrak M_2\mathfrak Q_1,
\end{equation}
where $\mathfrak M_j=D_q(f_j^+-f_j^-)$ and
$\mathfrak Q_j=t(f_j^++f_j^-)/(2D_q)$ at $a=1$.
The independent checker solves the Fredholm equations and compares actual
forward and reverse phases to this product, in both channels.
\checks{identities}\checks{scan}

\section{Correct half-line solution and a new reference correction}
\C Use only the corrected Wiener--Hopf normalization. Define
\begin{equation}
 g(s)=s^{-1/2}\Gamma(1+s/(2\pi))
       \exp\{-s\log[s/(2\pi e)]/(2\pi)\},\qquad
 \mathcal LG=\sqrt\pi\,g(s)/s^2.
\end{equation}
Put $G(u)=\frac43 y(u)^{3/2}$, $B=G'$, $f=G''$, $c=\sqrt\pi$.
Then $B(0^+)=c$, $\mathcal LB=\sqrt\pi g/s$ and
$\mathcal Lf=\sqrt\pi(g-1)$ for the ordinary derivative on $u>0$;
the removed constant is the contact delta. The dressed Jacobian reads
$B(u(y))=2\sqrt y/u'(y)$.

\C For $u>v>0$ the corrected, audited forward convolution is
\begin{equation}
 H(u,v):=\lim_{h\to0}N_{12}/h
 =2\int_0^v f(u-w)B(v-w)\dd w.
 \label{eq:H}
\end{equation}
Equation \eqref{eq:Uexact} and the endpoint asymptotics
$f_j^+\to1-B(u_j)/c$, $f_j^-\to0$ give
\begin{align}
 U(v)&:=\lim U_{12}/h=2c[B(v)-c],\\
 \widehat H(u,v)&=H(u,v)-U(v),\\
 \widehat H(u,v)+\widehat H(v,u)&=2[B(u)-c][B(v)-c],\\
 \boxed{H(u,v)+H(v,u)=2B(u)B(v)-2\pi.}
 \label{eq:two_recip}
\end{align}
The $-2\pi$ must not be dropped. In particular $H(u,u)=B(u)^2-\pi$;
the four-term phase vanishes on either edge.

\C There is a further, independently checkable correction to the frame
audit. The directional convolution and the four-term function satisfy
\begin{align}
 (\partial_u+\partial_v)H&=2f(u)B(v),\\
 (\partial_u+\partial_v)\widehat H&=2f(u)B(v)-2c f(v).
 \label{eq:PDE}
\end{align}
The frame audit instead prints $2f(u)[B(v)-c]$ for the latter source.
Its difference from the correct source is $2c[f(u)-f(v)]$, generally nonzero.
This is a reference mismatch, not a new homogeneous freedom. Given the
correct source and axis data, the characteristic equation still fixes the
phase. The exact endpoint identity, finite-sea tests, symbolic bookkeeping
and differentiated convolution are used, rather than the incorrect printed
four-term PDE. No conclusion that arbitrary two-body additions are allowed
is drawn from this correction.
At $(u,v)=(4,1)$, independent fourth-order differentiation of the
convolution gives $-0.347104254662920476$, agreeing with the corrected source
to $2.07\,10^{-17}$; the printed source is $+0.837748060692060127$ and
differs by $1.18485231535498060$.
\checks{identities}\checks{edge}\checks{analysis}

\section{Extracting the one-loop coefficient at fixed dressed momentum}
\C Denote $y_1=Y$, $y_2=r^2Y$, $0<r<1$. First take $h\to0$ at fixed
$(Y,r)$, then $Y\to\infty$. Each finite-$h$ momentum is exactly inverted;
the classical edge phase is
\begin{equation}
 H_{\rm cl}=4rY+2Y(1-r^2)\log\frac{1-r}{1+r}.
\end{equation}
At finite $Y$, $H-H_{\rm cl}$ contains higher inverse-$Y$ corrections;
it is not by itself a pure $O(\hbar^0)$ coefficient.

\C The corrected momentum expansion has
\begin{equation}
 u(y)=y-A\log y+c_{\rm edge}+O(y^{-1}),\qquad A=\frac1{2\pi},
 \qquad B(u(y))=2\sqrt y\,[1+A/y+O(y^{-2})].
\end{equation}
The constant $c_{\rm edge}$ cancels from the leading two-body correction.
To see the cancellation, substitute $w=v-u(z)$ in \eqref{eq:H}:
\begin{equation}
 H=4\int_0^{r^2Y}\sqrt z\,
 f\bigl(u(Y)-u(r^2Y)+u(z)\bigr)\dd z.
 \label{eq:regular}
\end{equation}
With $S=Y-r^2Y+z$, expansion of the integrand at fixed positive $z/Y$
gives $f=S^{-1/2}-\frac A2 S^{-3/2}\log[S r^2/z]+\cdots$.
Thus the leading residual reduces to
\begin{align}
 R_{12}^{(1)}
 &=-2A\int_0^{r^2}\frac{\sqrt t}{(1-r^2+t)^{3/2}}
       \log\frac{(1-r^2+t)r^2}{t}\dd t\\
 &=8A\int_0^r\log(s/r)\frac{s^2}{1-s^2}\dd s
 =-8A[\chi_2(r)-r].
 \label{eq:c0}
\end{align}
The change of variable is $s=\sqrt{t/(1-r^2+t)}$.
The integrals and their equality to the dilogarithms are checked at 70 digits.
\D The expansion at $z=O(1)$ requires a matched endpoint estimate to turn
this into a uniform theorem for the finite-sea equation. Direct inversion
and unconstrained fits test the result without assuming its intercept.

\C From \eqref{eq:two_recip}, subtracting the classical sum $8rY$ yields
\begin{equation}
 S^{(1)}=4A(r+r^{-1})-\pi,\qquad
 \boxed{\vartheta_N^{(1)}=\pi-8A\chi_2(r)+4A(r-r^{-1}).}
 \label{eq:answer}
\end{equation}
For exchanged arguments the antisymmetric coefficient changes sign.
At $r=1$ the finite-$Y$ antisymmetric phase is zero, whereas the ordered
large-$Y$ limit of \eqref{eq:answer} is $\pi/2$ as $r\uparrow1$.
The near-equal-amplitude boundary layer is therefore not covered by this
ordered expansion. The $r\downarrow0$ limit is also nonuniform: the forward
coefficient tends to zero while its symmetric and antisymmetric parts have
opposite $1/r$ terms. They must be combined before taking a soft-soliton limit.

\C The one-body calculation supplies a nontrivial consistency check:
\begin{equation}
 C=s_1-A/2,\qquad s_1=3/(4\pi),\qquad
 C=1/(2\pi),\qquad A=2(s_1-C).
\end{equation}
Consequently the coefficient of $r+r^{-1}$ in the symmetric phase is
$8(s_1-C)$, agreeing with the independently inverted Jacobian. A common
sound-line shift changes $C$ and $s_1$ equally and leaves this coefficient
unchanged. The decomposition found here is an edge dilogarithm, a rational
Jacobian term, and an endpoint constant. No additional independent
sound-speed two-body term is found in this strict corner. Splitting the
rational term into separate multiples of $C$ and $s_1$ is algebraically
possible but is not an independently established separation of mechanisms.
\checks{identities}\checks{edge}\checks{analysis}

\section{Head-on channel and why reference data must accompany the answer}
\C For opposite edges, $f_2^+=O(h)$ while $f_2^-\to B(v)/c-1$.
The opposite-edge directional interaction and its edge reference are
$N_{12}/h=O(h)$ and $U_{12}/h=O(h)$ at fixed $y_j$.
The leading classical terms are $\delta_T/h\to4\sqrt{y_1y_2}$ and
$C_{12}/h\to-4\sqrt{y_1y_2}$. Thus the leading directional residual,
its symmetric part and its antisymmetric part all tend to zero.
New finite-sea extrapolations test both argument orders. This assertion
concerns the stated iterated corner, not finite-amplitude head-on scattering
or a global crossing continuation of \eqref{eq:answer}.

\C The algebraically antisymmetric endpoint representative in the audited
frame dictionary is
$\mathcal B_{12}=\widehat\Phi_{12}-\mathfrak M_1\mathfrak Q_2$.
It is unique once this particular subtraction is specified. Its edge values
are, however,
\begin{align}
 \mathcal B_{\rm ov}/h
 &=H-B_1B_2+c(B_1-B_2)+\pi,\\
 \mathcal B_{\rm ho}/h&=(B_1-c)(B_2-c).
 \label{eq:Bchannels}
\end{align}
Relative to their classical phases, these contain respectively
\begin{align}
 \mathcal B_{\rm ov}/h-\delta_{{\rm cl},\rm ov}
 &=2\sqrt\pi\sqrt Y(1-r)+\vartheta_N^{(1)}+O(Y^{-1/2}),\\
 \mathcal B_{\rm ho}/h-4rY
 &=-2\sqrt\pi\sqrt Y(1+r)+\pi+\frac2\pi(r+r^{-1})+O(Y^{-1/2}).
 \label{eq:Btails}
\end{align}
At fixed small $\beta$, $\sqrt Y=\beta_1/\sqrt{\hbar_B}$.
Discarding these terms silently would change the quantum expansion and, in
the head-on channel, the reported constant. The JSON retains $U_{12}$,
$U_{21}$, the four-term phase and $\mathcal B_{12}/h$ as well as $R$.

\D Removing an arbitrary antisymmetric addition from the Toda Fredholm
problem does not by itself identify any of these representatives with the
continuum canonical scattering operator. The frame audit explicitly leaves
that endpoint-operator identification open, and the Phase 2i audit does not
select an absolute phase. No new choice of regulator, physical frame,
statistics offset or crossing branch is made here. Accordingly the table
for $\vartheta_N$ must not be relabelled as the continuum Gamma-relative
$\vartheta$ without this additional map.
\checks{identities}\checks{scan}\checks{analysis}

\section{Numerical extraction and error budget}
\C The finite-sea mesh uses eleven values
\[
h\in\{.08,.06,.04,.03,.02,.015,.01,.0075,.005,.00375,.0025\},
\]
with $N=\lceil10/h+200\rceil$, plus 25\% refinements at $.01,.005,.0025$.
The maximum mesh has 5250 points. Its Gauss weights are recomputed from
long-double Newton iteration and a Legendre recurrence after the initial
SciPy-weight run failed reciprocity; the sea linear algebra remains float64.
Polynomial quadrature moments are checked independently. Fifteen pairs
$Y=4,8,16,32,64$, $r=.4,.5,.6$ are solved in both channels and both argument
orders, with exact endpoint $q$ and individually inverted dressed momenta.
Every solve checks counting, momentum inversion, endpoint identities and
reciprocity. The $h$ fits are evaluated at 60 digits; their float64 inputs
are not represented as high-precision sea data.

\C The central extrapolation is quartic on the last seven $h$ points.
Alternatives are cubic on six, quintic on eight, quadratic on four, and
cubic plus $h\log h$ or $h^2\log h$ on seven points. No target intercept
is inserted. The $h$ uncertainty is the largest change of intercept.
The mesh error is the largest refined-minus-base difference plus
$2\,10^{-10}$, propagated through the largest extrapolation weight norm.
The two errors are added. The same method separately checks the vacuum
sound coefficient against $3/(4\pi)$.

\C The half-line grid has 36 pairs: nine scales from $4$ to $1024$ at
$r=.4,.5,.6$, and scales $32,128,512$ at $r=.25,.75,.9$.
Each convolution uses 32-digit Talbot inversion and geometrically segmented
16- and 24-point mpmath Gauss rules, resolving the $u=O(1)$ layer.
De Hoog inversion at 45 digits independently checks four one-body points.
The largest cross-method difference is $2.57\,10^{-30}$ and the largest
16/24-node convolution difference is $1.01\,10^{-15}$.
The derivative entering the convolution is separately checked by centered
differences of the inverted $B$. Quadrature differences are propagated
through the amplitude-fit weights. These diagnostics do not bound a common
systematic error of the half-line model.

\C For the three dense ratio grids, the central large-$Y$ fit is cubic
in $1/Y$ through the last four points. Alternatives include quadratic
fits, a $Y^{-3/2}$ term, a $Y^{-2}\log Y$ term, and a two-point linear
fit. For the three sparse ratio grids the quadratic intercept and the
last-two linear intercept are compared. All intercepts are free. Their
largest spread is the amplitude-tail error; it is added to the propagated
quadrature diagnostic. A separate table compares the finite-sea intercepts
against the half-line model rather than merging their errors into a
fictitiously rigorous bound on the limiting physical observable.

\C The following numbers are generated by \file{twobody_toda_analysis.py} from the new caches.
\begin{center}\small\begin{tabular}{rrrrr}\toprule
$r$ & Forward intercept & Tail + quadrature & $\vartheta_N$ intercept & Tail + quadrature\\\midrule
$0.25$ & $-0.00226183803$ & $1.79\,10^{-7}$ & $0.4336115256$ & $0.00033308$\\
$0.40$ & $-0.00962285926$ & $7.65\,10^{-9}$ & $1.2857724554$ & $3.09\,10^{-6}$\\
$0.50$ & $-0.01951540939$ & $6.88\,10^{-10}$ & $1.5305278134$ & $1.13\,10^{-6}$\\
$0.60$ & $-0.03545557168$ & $9.61\,10^{-9}$ & $1.6631322646$ & $5.30\,10^{-7}$\\
$0.75$ & $-0.07717786602$ & $3.83\,10^{-7}$ & $1.7381235642$ & $2.23\,10^{-6}$\\
$0.90$ & $-0.16035085391$ & $2.04\,10^{-6}$ & $1.7009286895$ & $3.11\,10^{-6}$\\
\bottomrule\end{tabular}\end{center}
\C For $r=1/2$, the proposed exact antisymmetric value is $1.53052781327896308\ldots$. The fit did not use this value.
\begin{center}\small\begin{tabular}{rrrrr}\toprule
$Y$, $r=1/2$ & Sea $R_{12}$ & Mesh error & $h$-fit spread & Difference from WH\\\midrule
$4$ & $-0.0248313244$ & $8.01\,10^{-9}$ & $3.11\,10^{-8}$ & $2.62\,10^{-11}$\\
$16$ & $-0.0208546047$ & $9.47\,10^{-9}$ & $1.70\,10^{-6}$ & $7.09\,10^{-9}$\\
$64$ & $-0.0198545583$ & $1.53\,10^{-8}$ & $0.00011925$ & $4.16\,10^{-6}$\\
\bottomrule\end{tabular}\end{center}
\begin{center}\small\begin{tabular}{rrrrr}\toprule
$Y$, $r=1/2$ & Head-on $R_{12}$ & Total diagnostic & Head-on $R_{21}$ & Total diagnostic\\\midrule
$4$ & $-6.37\,10^{-11}$ & $3.75\,10^{-8}$ & $1.16\,10^{-10}$ & $1.66\,10^{-7}$\\
$16$ & $-4.55\,10^{-8}$ & $2.16\,10^{-6}$ & $7.13\,10^{-8}$ & $7.71\,10^{-6}$\\
$64$ & $-1.92\,10^{-5}$ & $0.00019692$ & $3.35\,10^{-5}$ & $0.00062092$\\
\bottomrule\end{tabular}\end{center}
\C Restricted three-parameter candidate fits to all six independent numerical antisymmetric intercepts give:
\begin{center}\small\begin{tabular}{lr}\toprule
Basis & Maximum absolute residual\\\midrule
$1,\chi_2(r),r-r^{-1}$ & $1.44\,10^{-5}$\\
$1,r,r^{-1}$ & $0.01579042$\\
$1,\log r,\log[(1-r)/(1+r)]$ & $0.08669543$\\
$1,r-r^{-1},\log r\log[(1-r)/(1+r)]$ & $0.02143624$\\
\bottomrule\end{tabular}\end{center}
\C The freely fitted coefficients in the first row are $(3.14175645, -1.27339438, 0.63667460)$, compared with $(\pi,-4/\pi,2/\pi)$. This tests the specified finite bases; it is not an exclusion theorem for all rational or logarithmic functions.
\C This unweighted fit is affected by the less precise $r=0.25$ intercept. The proposed exact formula is tested separately against each point and its own uncertainty, and all six comparisons pass; the unweighted residual is not its error budget.
\C The independently extrapolated sound coefficient is $0.238732414641\pm 8.99\,10^{-9}$, compared with $3/(4\pi)=0.238732414637843\ldots$.
\C Recorded times for completed dense-solve and convolution rows are 953 s and 1136 s respectively. These heavy jobs were run concurrently; the default rerun validates their caches. The failed refinement and subsequent identity checks are not included in these cached-row totals.


\D The quoted errors are deterministic sensitivity envelopes, not
statistical confidence intervals and not a proof that no other asymptotic
term is present. In particular, a small contour-inversion difference does
not bound finite-sea matching, and no sampled estimate controls $r=0$ or
$r=1$. Full quantum equivalence is not inferred from agreement of this
coefficient with a possible future continuum result.
\checks{scan}\checks{edge}\checks{analysis}

\section{Verification summary, failures and reproducibility}
\begin{longtable}{@{}L{5.4cm}L{2.1cm}L{7.3cm}@{}}
\toprule Claim & Status & Check and limitation\\\midrule\endhead
Endpoint evaluation, reciprocity and dressed Jacobian & (c) &
\file{twobody_toda_identities.py}; independent 70-digit quadrature and root derivative.\\
Correctly referenced edge phase and reverse phase & conditional (c) &
\file{twobody_toda_edge.py}, \file{twobody_toda_scan.py}; two-term/four-term terms retained.\\
Forward dilogarithmic coefficient & conditional (c) &
Independent integral, direct inverse-Laplace convolution, free intercepts.\\
Symmetric coefficient and one-body consistency & (c) in edge model &
Exact reciprocity and same $A=2(s_1-C)$, tested numerically.\\
Antisymmetric directional function and head-on zero & conditional (c) &
Both finite-sea orders and channels; explicit normalization in JSON.\\
Printed frame-audit four-term PDE & refuted (c) &
Endpoint reference differentiation; discrepancy $2c[f(u)-f(v)]$.\\
Physical scalar phase and Gamma-relative residual & (d) &
Endpoint operator map, absolute phase and crossing normalization remain open.\\
Uniform fixed-nonzero-$P$ limit & (d) &
No uniform finite-sea/large-$Y$ or near-diagonal estimate supplied.\\
\bottomrule
\end{longtable}

\paragraph{Failures recorded.}
\begin{enumerate}
\item The first finite-sea run failed before producing a row because the
node-count attribute \texttt{N} shadowed a method of the same name.
The method was renamed \texttt{directional}; no equation, grid or tolerance
was changed.
\item The first 5250-node refinement failed the unchanged $2\,10^{-8}$
reciprocity tolerance. The original assertion did not print the residual,
so its exact size was not retained; the failed row was not accepted into
the cache. Recomputing the high-order Gauss weights using extended-precision
Legendre recurrence addresses the quadrature accuracy rather than relaxing
the test. The replacement changes the largest weight by $1.40\,10^{-13}$
and passes with maximum reciprocity residual $2.58\,10^{-11}$ and counting
error $1.78\,10^{-15}$. The successful replacement is stored in the cache.
\item The printed four-term Krein source in the frame audit fails the exact
reference differentiation. The wrong source is not used, and the discrepancy
is retained as a positive refutation check. This does not weaken the
Fredholm uniqueness condition.
\item Restricted rational-only and logarithm-only candidate fits are tested
against the numerical intercepts; their nonzero residuals are reported in
the table and JSON. Their rejection is limited to these tested bases.
\item The PDF skill's artifact-marker helper was not installed in this local
workspace (\texttt{MODULE\_NOT\_FOUND}); the authorized local LaTeX build
is used directly. An exploratory read also requested a nonexistent historical
solver filename; the actual solver sources were then located. Neither
operation modified an existing file.
\item The first PDF build had a 5.24-point overfull line containing the
entire $h$ grid; moving the grid to display math repaired it. Nonbreaking
spaces also keep claim labels with their following text at page boundaries.
\end{enumerate}
No mathematical check was weakened to hide a failure. The success markers
certify the asserted checks, not the open physical identifications.
All four final scripts printed \texttt{ALL CHECKS PASSED}.

\paragraph{Reproduction and interface.}
Run only the new scripts, in this order:
\begin{verbatim}
python3 phase2/scripts/twobody_toda_identities.py
python3 phase2/scripts/twobody_toda_scan.py
python3 phase2/scripts/twobody_toda_edge.py
python3 phase2/scripts/twobody_toda_analysis.py
\end{verbatim}
The two heavy scripts resume their own JSON caches; \texttt{--refresh}
regenerates them. The coordinator-facing \file{twobody_toda_grid.json}
contains definitions, signed channels, raw finite-$h$ fit variants, the
finite-$Y$ half-line grid, the limiting coefficients, candidate residuals
and separate numerical error components. The raw sea and edge caches are
also preserved. Build the PDF by running \texttt{pdflatex} twice from
\file{phase2/}; the numerical table fragment is supplied alongside this source.

\paragraph{Internal sources read; no new external bibliography.}
\begingroup\raggedright
The sources used are \file{AGENTS.md}, \file{README.md}, the conclusions
and relevant Toda sections of \file{phase2/report.tex},
\file{frame_dictionary.tex}, \file{audit_frame_dictionary.tex},
\file{audit_phase2i.tex}, \file{onebody_toda_precision.tex},
\file{audit_onebody_scheme.tex}, the residual-phase sections of
\file{smx_candidates.tex} and \file{audit_smx.tex}, and the relevant
\file{errata.tex} and \file{audit_toda_dictionary.tex} corrections.
Implementation sources were inspected for conventions and methods;
none of their executable modules was imported or run.
All numerical results in this note are new task-prefixed outputs.
\par\endgroup
\end{document}
